\section{Background}\label{sec:background}

The framework built on four established components, each contributing a distinct piece of the synthetic option generation pipeline: the JumpHMM produced realistic multi-asset price paths; the Heston stochastic variance process provided a mechanism for converting price dynamics into implied volatility; a recombining binomial lattice translated IV into American option prices with early exercise; and a neural surrogate replaced the analytic shape function for the IV surface with a data-driven representation that absorbed sector-specific smile and term-structure geometry.

\subsection{Jump Hidden Markov Model}

The Jump Hidden Markov Model (JumpHMM)~\cite{jumphmm2025} was developed for generating synthetic equity time series that simultaneously preserved three stylized facts of financial returns: heavy tails, negligible linear autocorrelation, and persistent volatility clustering~\cite{cont2001}. The model discretized continuous excess growth rates into $N$ states via equal-probability quantile bins of a fitted Laplace distribution. The hidden state $S_t \in \{1, \ldots, N\}$ at each timestep governed the distribution of the excess growth rate $G_t$; each state $k$ emitted observations from a location-scale Student-$t$ distribution: $G_t \mid S_t\!=\!k \;\sim\; \mu_k + \sigma_k \cdot t_\nu$, where $\mu_k$ and $\sigma_k$ were learned per-state location and scale parameters and $\nu$ was the shared degrees-of-freedom parameter. At each timestep, with probability $\epsilon$ a Poisson$(\lambda)$-distributed jump forced the chain into tail states (bottom or top $N_{\text{tail}}$ states), creating sustained dwell periods in extreme-volatility regimes that produced the volatility clustering characteristic of empirical return series. For multi-asset generation, cross-asset dependence was imposed via a Student-$t$ copula~\cite{demarta2005} that preserved each asset's marginal distribution while introducing realistic tail dependence; this copula structure ensured that joint tail events, in which multiple assets simultaneously occupied extreme states, occurred more frequently than a Gaussian copula would imply.

\subsection{Heston stochastic variance model}

The Heston model~\cite{heston1993} specified the instantaneous return variance $v$ of the underlying asset (i.e., the variance rate of its log-returns) as a mean-reverting square-root process:
\begin{equation}\label{eq:heston_standard}
    dv = \kappa(\theta - v)\,dt + \sigma_v \sqrt{v}\,dW_v
\end{equation}
where $\kappa$ was the mean-reversion speed controlling how quickly the variance returned to its long-run level, $\theta$ was the long-run variance target, and $\sigma_v$ was the vol-of-vol governing the amplitude of stochastic fluctuations in $v$. In the standard formulation, the asset price and variance processes were coupled through a correlation $\rho$ between their driving Brownian motions, producing the leverage effect in which negative returns coincided with rising volatility. The key property for our purposes was mean-reversion: the variance process was pulled toward $\theta$, so by controlling $\theta$ we could control the equilibrium level of implied volatility. More generally, the Heston model produced a full European IV surface in $(K,\tau)$ through its semi-closed-form characteristic-function pricing, and this analytic surface was the shape function that the neural surrogate below replaced in our pipeline. In the standard model $\theta$ was a constant; our extension made $\theta$ time-varying and state-dependent, linking it to the JumpHMM regime state, contract characteristics, and market conditions.

\subsection{Lattice models for American options}

Once the Heston shape function above supplies an IV per contract, the pipeline still needs to convert $(S, K, \tau, r, q, \sigma)$ into an option price; for American options this rules out closed-form Black-Scholes and requires a numerical scheme. Lattice methods discretize the underlying-price process onto a recombining tree and price contingent claims by backward induction with an early-exercise check at each node~\cite{broadie1996}. The Cox-Ross-Rubinstein (CRR) tree~\cite{cox1979} is the canonical construction, with up/down factors $u = e^{\sigma\sqrt{\Delta t}}$, $d = 1/u$, and risk-neutral probability $p = (e^{r\Delta t} - d)/(u - d)$. Subsequent variants modified the moment-matching to address specific shortcomings of the CRR construction: Jarrow-Rudd~\cite{jarrow1983} matched the first two moments under the risk-neutral measure with $p = 1/2$ at the cost of a non-zero log-drift term, Tian's modified lattice~\cite{tian1993} added a third-moment match to suppress the order-$\sqrt{\Delta t}$ pricing bias of the CRR tree, and Leisen-Reimer (LR)~\cite{leisen1996} replaced the moment-matching construction altogether with a Peizer-Pratt inversion of the binomial CDF, choosing $p$ and $u$ at each step so that the strike $K$ landed on a tree node by design at every $(S, \sigma)$ configuration. All four constructions converge to the Black-Scholes price as the number of steps increases, but their finite-step behavior differs sharply: the CRR price oscillates around the continuum limit because $K$ drifts across node boundaries as $\Delta t$ shrinks, the LR price converges monotonically because the strike is pinned to a node, and the Tian and Jarrow-Rudd trees fall in between~\cite{leisen1996,broadie1996}. The pinned-strike property of the LR construction also stabilizes finite-difference Greeks: small bumps to $S$ or $\sigma$ leave the strike on a node, so the staircase aliasing that CRR exhibits when the bump moves $K$ across a node boundary is absent.

Two of these lattices appear later in this work and the choice between them is dictated by the use case. The pricing-pipeline pricer is the CRR tree at $200$ steps: with IV pre-computed by the upstream Heston shape function, the pricer needs only a single $\sigma$ input per contract, and CRR's recombining structure makes the cost of pricing the full $234{,}549$-row ladder linear in the number of steps. The forward-scenarios pricer, in contrast, required smooth Greeks at modest tree depth because the per-path scenarios in the §5 short-premium analysis computed $\Delta$, $\Gamma$, and Vega by central finite differences against thousands of $(S_t, \sigma_t)$ configurations along each simulated path; we therefore switched to the LR tree at $201$ steps for that purpose, where the pinned-strike construction eliminated the bump-induced aliasing that CRR produced at the same depth. The pipeline did not depend on either choice in particular: any of the four lattices above would have been a drop-in replacement for the appropriate stage, and the decomposition between IV calibration and option pricing kept the choice of tree orthogonal to the rest of the framework.

\subsection{Neural representations of the implied volatility surface}

Parametric IV models such as SABR~\cite{hagan2002} and SVI~\cite{gatheral2004} expressed the smile and term structure through low-dimensional closed-form functions, chosen to interpolate observed quotes while respecting no-arbitrage constraints. These forms were parsimonious and interpretable but inherited their expressive capacity from the chosen functional family, and a single global parameterization across a broad universe of tickers could not simultaneously accommodate sector-specific skew patterns and smile curvatures. A complementary line of work~\cite{horvath2021,bayer2019} replaced the analytic form with a neural network that took $(\ln\tau,\ln m)$ or a similar low-dimensional set of contract and model features as input and returned an IV or variance output; the network provided a flexible universal approximator for the shape of the surface while the downstream pricing pipeline remained unchanged. Neural surrogates of this kind have been used both as fast pricing approximators and as calibration targets in deep calibration schemes, and they admit the same multiplicative decomposition between a per-ticker level and a shared shape that the parametric forms use, making them a drop-in replacement for the functional form inside a generative pipeline. We used the neural surrogate in the structural sense rather than the deep-calibration sense: $\psi_{\mathrm{NN}}$ was the shape factor inside the multiplicative decomposition whose level $v_t$ evolved with a JumpHMM-coupled Heston process, so the network supplied the cross-sectional geometry while the dynamics came from the generative model.
